\documentclass[11pt]{article}
\usepackage[margin=0.82in]{geometry}
\usepackage{amsmath,amssymb}
\usepackage{booktabs}
\usepackage{graphicx}
\usepackage{float}
\usepackage{placeins}
\usepackage[hidelinks]{hyperref}
\usepackage{xcolor}

\graphicspath{{../outputs/figures/}}
\setlength{\parindent}{0pt}
\setlength{\parskip}{0.55em}
\setlength{\tabcolsep}{5pt}
\newcommand{\todo}[1]{\textcolor{red}{\textbf{TODO: #1}}}
\newcommand{\E}{\mathbb{E}}
\newcommand{\R}{\mathbb{R}}

\title{ME2501 --- Assignment 1\\Exploring Dictionary Learning}
\author{Haina Wang, 260010205 }
\date{Fall 2026}

\begin{document}
\maketitle

\paragraph{Repository.}
\todo{Insert the URL of the private repository in the
\texttt{hku-data8014-students} organization.}

\section{Implementation and experimental protocol}

I added a NumPy/SciPy package, \texttt{data8014/dictlearn}, containing the
synthetic generators, recovery and validation metrics, ISTA/FISTA, MSP,
alternating minimization, convergence fitting, and MNIST PCA/whitening.  Thin
experiment drivers in \texttt{scripts/dl\_part*.py} write incremental JSONL
logs and regenerate figures from those logs.  Part 3 uses \texttt{float32};
the synthetic experiments use \texttt{float64}.  Unit tests check moments and
dictionary constraints, sign/permutation invariance, the exact one-step ISTA
solution for an orthogonal dictionary, step-size calculations, and the Part 4
variance construction.

Unless otherwise stated, reported uncertainty is the sample standard deviation
across independent seeds.  I used five seeds in Parts 0--2 and 4 and three
seeds in the more expensive MNIST experiments.  For recovery I primarily use
(3), report duplicate matches, and additionally log a Hungarian one-to-one
score.  Thus a collapsed dictionary cannot silently appear as successful.
Validation codes are always inferred from held-out samples.  In the complete
case, because every MSP iterate is orthogonal, ISTA from zero gives the exact
LASSO solution in one step,
\[
  \widehat Z=\operatorname{soft}(X_{\rm val}\widehat A^\top,\lambda/2).
\]

My MSP update is
\[
 A_{t+1}=\mathcal P_{O(D)}\!\left(
       ((XA_t^\top)^{\circ 3})^\top X\right),
\]
where the polar projection is computed by an SVD.  One alternating-minimization
block consists of 30 warm-started ISTA steps, one dictionary-gradient step,
and row normalization.  I recompute
$\eta=(2\|A\|_2^2)^{-1}$ and $\nu=(2\|Z\|_2^2)^{-1}$ in every block.

\section{Part 0: setup and warm-ups}

\subsection{Generator, metrics, and MSP sanity check}

For the Bernoulli--Gaussian generator I draw
$z_k=b_kg_k$, $b_k\sim\mathrm{Bern}(\theta)$ and
$g_k\sim\mathcal N(0,1/\theta)$.  Complete dictionaries are QR factors;
overcomplete dictionaries have independently drawn, unit-normalized rows.
With the prescribed MSP sanity setting, all five seeds recovered every atom to
$s_k>0.99$ in 6--8 one-based iterations.  Final minimum correlations ranged
from 0.9986 to 0.9993.

\begin{figure}[H]
  \centering
  \includegraphics[width=0.82\linewidth]{dl_part0_sanity_msp.png}
  \caption{MSP sanity check.  The minimum atom correlation exceeds 0.99 well
  before the 30-iteration reference budget for every seed.}
\end{figure}

\subsection{ISTA convergence}

For the fixed ground-truth overcomplete dictionary, the smooth term has
Lipschitz constant $L=2\|A\|_2^2$.  The plotted worst-case proximal-gradient
bound is
\[
 f(Z_k)-f^\star\leq \frac{L\|Z_0-Z^\star\|_F^2}{2k}.
\]
I estimated $f^\star=8814.0855168404$ by a long FISTA run.  The solution has
8273 nonzero entries out of 64000.  ISTA permanently identifies this entrywise
support at iteration 2273; after that, the fitted geometric factor is
$\rho=0.99783$ per iteration.  The eventual geometric behavior is faster than
the generic $O(1/k)$ guarantee.  The bound does not use the finite active-set
structure: after signs and support stabilize, soft thresholding is affine and
the remaining iteration is a linear contraction on the active coordinates.

\begin{figure}[H]
  \centering
  \includegraphics[width=0.95\linewidth]{dl_part0_warmup_ista.png}
  \caption{ISTA objective gap against the worst-case convex bound (left) and
  the geometric tail after support identification (right).}
\end{figure}

\subsection{Why orthogonality costs no generality}

Let a noiseless complete model have an arbitrary invertible $A$ and
$x=A^\top z$.  Since the coordinates of $z$ have unit variance,
\[
 \Sigma_x=\E[xx^\top]=A^\top A.
\]
With the symmetric whitening matrix $W=\Sigma_x^{-1/2}$, transform row samples
as $\bar X=XW$ and define $\bar A=AW$.  Then
\[
 \bar A^\top\bar A=W^\top A^\top A W=I.
\]
Because $\bar A$ is square, it is orthogonal.  Learning the whitened model
$\bar X=Z\bar A$ therefore loses no information: the original dictionary is
$A=\bar A W^{-1}$.  In finite samples I use the empirical covariance; this
argument requires the stipulated invertibility and zero-noise setting.

\section{Part 1: convergence and sample-size scaling}

\subsection{Optimization convergence (T-axis)}

I used $D=16$, $N=8000$, $\theta=0.1$, $\sigma=0$, $\lambda=0.6$, and five
seeds.  Both metrics reach their finite-sample/LASSO plateaus near $T=12$.
To avoid fitting the plateau as if it were decay, I estimated a tail floor and
fit only the positive transient above it.  Table~\ref{tab:p1t} compares the
mean-curve fits; the per-seed parameter spread is also shown.

\begin{table}[H]
\centering
\caption{Part 1(a) floor-aware transient fits.  Larger $R^2$ is better.}
\label{tab:p1t}
\begin{tabular}{lrrrrrr}
\toprule
Metric & floor & power $\alpha$ & $R^2$ & exp. $\rho$ & $R^2$ & per-seed $\rho$ \\
\midrule
Recovery error & $3.61\!\times\!10^{-4}$ & 4.307 & .592 & .350 & .799 & $.210\pm.126$ \\
Eval loss      & .1470 & 4.957 & .608 & .324 & .814 & $.257\pm.098$ \\
\bottomrule
\end{tabular}
\end{table}

The exponential model is more consistent with both measured transients, but
$R^2\approx0.8$ and the short, very rapid approach to a floor mean that neither
simple law is a complete global description.  Recovery tends to zero up to
finite-sample error, whereas the eval loss retains a nonzero LASSO shrinkage
floor, so the raw metrics do not share the same asymptote.

\begin{figure}[H]
  \centering
  \includegraphics[width=0.93\linewidth]{dl_part1a_msp_convergence.png}
  \caption{Part 1(a): held-out and recovery errors on semilogy and log--log
  axes.  Bands show one standard deviation over five seeds.}
\end{figure}

\subsection{Sample-size scaling (N-axis)}

For each seed I drew the $N=10000$ dataset once and used prefixes, holding the
true dictionary, initialization, and validation data fixed.  I stopped using a
signed/permuted iterate-stationarity criterion rather than ground-truth
recovery.  All 45 runs converged before the cap (maximum 126 iterations).
The first six of nine $N$ values were fit and $N=10000$ was predicted.

\begin{table}[H]
\centering
\caption{Part 1(b) fits of $\mathcal E-\mathcal E_\infty=cN^{-\beta}$.}
\begin{tabular}{lrrrr}
\toprule
Metric & mean-curve $\beta$ & $R^2$ & per-seed $\beta$ & target relative error \\
\midrule
Recovery error & 1.285 & .936 & $1.212\pm.192$ & 16.1\% \\
Eval excess    & .472 & .907 & $.463\pm.077$ & 207\% \\
\bottomrule
\end{tabular}
\end{table}

The raw eval-loss prediction error is only 13.9\%, but this is misleading
because the estimated floor $\mathcal E_\infty=0.13673$ dominates.  Relative to
the quantity actually fit, $\mathcal E-\mathcal E_\infty$, the extrapolation
error is 207\%.  Thus recovery is reasonably compatible with a power law over
the measured range, while eval-excess extrapolation is not stable enough to
support the same claim despite its in-range $R^2$.

\begin{figure}[H]
  \centering
  \includegraphics[width=0.9\linewidth]{dl_part1b_msp_sample_size.png}
  \caption{Part 1(b): converged recovery error and held-out excess loss versus
  sample count.}
\end{figure}
\FloatBarrier

\section{Part 2: density and identifiability}

\subsection{Cost and transition}

For $\theta\in\{.2,.4,.6,.7,.8,.85\}$, the mean $T_\star$ values were
$8.0$, $9.8$, $17.8$, $26.6$, $48.0$, and $88.0$; their standard deviations
were $2.55$, $1.79$, $1.30$, $8.17$, $4.85$, and $16.51$, respectively.  No
run was censored.  After recording the first-passage time, I continued each
run to iterate stationarity before evaluating its held-out loss.  Density
therefore causes a large and increasingly variable optimization cost.

\begin{figure}[H]
  \centering
  \includegraphics[width=0.72\linewidth]{dl_part2_Tstar_vs_theta.png}
  \caption{Iterations required to recover at least 95\% of atoms above 0.9.}
\end{figure}

At higher density the mean recovered fractions were
\[
\begin{array}{c|rrrrrrr}
\theta & .88 & .90 & .91 & .92 & .93 & .95 & 1.00\\ \hline
\text{fraction} & .9875 & .9375 & .4875 & .175 & .075 & 0 & 0\\
\text{std.} & .028 & .088 & .277 & .103 & .052 & 0 & 0
\end{array}
\]
All runs reached iterate stationarity before the 20000-iteration cap.  The
sharp fall between .90 and .92 is evidence of a finite-sample/algorithmic
phase transition.  In contrast, converged held-out loss rises smoothly: its
mean is .2815 at $\theta=.2$, 1.1709 at .90, 1.1935 at .91, 1.2043 at .92,
and 1.2118 at 1.  There is no loss discontinuity corresponding to the recovery
collapse, so reconstruction is not an identifiability diagnostic.

\begin{figure}[H]
  \centering
  \includegraphics[width=0.95\linewidth]{dl_part2_transition.png}
  \caption{Part 2 transition.  Recovery changes sharply while held-out
  reconstruction loss changes smoothly.}
\end{figure}

\subsection{Why the transition occurs}

For $z=bg$, $b\sim\mathrm{Bern}(\theta)$ and
$g\sim\mathcal N(0,1/\theta)$,
\[
 \E[z^2]=\theta\frac1\theta=1,\qquad
 \E[z^4]=\theta\,3\left(\frac1\theta\right)^2=\frac3\theta,
 \qquad \kappa(\theta)=3\left(\frac1\theta-1\right).
\]
At the observed transition center $\theta=.91$, $\kappa\approx.297$; at
$\theta=1$, $\kappa=0$.

At $\theta=1$, $z\sim\mathcal N(0,I)$.  For every orthogonal $R$, set
$Z'=ZR^\top$ and $A'=RA$.  Then $Z'A'=ZA$, while rotational invariance gives
$Z'\stackrel{d}{=}Z$.  Infinitely many dictionaries therefore induce the same
data law.  MSP's fourth-power objective exploits coordinate-wise non-Gaussian
kurtosis: sparse directions produce unusually large fourth moments when the
iterate aligns with the independent components.  As $\theta\to1$, the excess
kurtosis and hence this directional contrast vanish.  The smooth loss remains
small because a rotated Gaussian representation reconstructs just as well.

\section{Part 3: MNIST}

\subsection{Spectrum, projection, and whitening}

I centered the MNIST training images, diagonalized their pixel covariance,
and used $k=32$.  The 80\%, 90\%, and 99\% variance thresholds require 44, 87,
and 331 components, respectively, so $k=32$ deliberately trades retained
variance for computational cost.  I divided each projected coordinate by the
square root of its training eigenvalue.  The resulting relative covariance
error was
$\|\widehat\Sigma_y-I\|_F/\|I\|_F=2.77\times10^{-7}$.

\begin{figure}[H]
  \centering
  \includegraphics[width=0.78\linewidth]{dl_part3_spectrum.png}
  \caption{MNIST covariance spectrum and cumulative explained variance.}
\end{figure}

\subsection{T-axis}

I trained on $N=12000$ images for 400 blocks with $\lambda=2$ and evaluated on
a fixed, unused set of 2000 test images.  Table~\ref{tab:p3t} reports the final
mean and standard deviation.

\begin{table}[H]
\centering
\caption{Part 3 T-axis results after 400 blocks (three seeds).}
\label{tab:p3t}
\begin{tabular}{lrrrr}
\toprule
$m$ & held-out loss & random loss & learned/random & code nnz fraction \\
\midrule
64  & $10.1889\pm.0098$ & $15.9189\pm.3085$ & $.6402\pm.0128$ & $.13050\pm.00084$ \\
128 & $7.9752\pm.0238$  & $11.8963\pm.0640$ & $.6704\pm.0043$ & $.07403\pm.00040$ \\
\bottomrule
\end{tabular}
\end{table}

Floor-aware held-out fits favored exponential decay: for $m=64$,
$\rho=.9789$ with $R^2=.954$ versus power $R^2=.841$; for $m=128$,
$\rho=.9850$ with $R^2=.822$ versus power $R^2=.774$.  Training-objective
exponential fits had $R^2=.938$ and .941.  Thus the real-data overcomplete
problem also has an approximately geometric transient, but it relaxes over
hundreds of blocks rather than MSP's roughly twelve iterations, and the late
curves are dominated by a floor and small nonmonotonic fluctuations.

\begin{figure}[H]
  \centering
  \includegraphics[width=0.94\linewidth]{dl_part3_taxis.png}
  \caption{MNIST T-axis, with one-standard-deviation bands and random-dictionary
  held-out baselines.}
\end{figure}

\subsection{N-axis}

For each seed I fixed an image ordering and used prefixes at the seven required
sample counts.  I fit the smallest five values with a free asymptotic floor and
predicted $N=12800$.

\begin{table}[H]
\centering
\caption{MNIST sample-size fits and diagnostics.}
\begin{tabular}{lrrrrr}
\toprule
$m$ & $\beta$ & per-seed $\beta$ & $R^2$ & prediction error & capped runs \\
\midrule
64  & .558 & $.555\pm.155$ & .9999 & 1.22\% & 3/21 \\
128 & .485 & $.494\pm.184$ & .9996 & 1.54\% & 12/21 \\
\bottomrule
\end{tabular}
\end{table}

The large $R^2$ and small target errors support a floor-plus-power description
on these points, but 15 of 42 runs reached the permitted 1000-block cap rather
than the $10^{-5}$ stopping criterion.  This is especially important for
$m=128$ and limits precision.  The synthetic eval exponent was
$.463\pm.077$ per seed.  The $m=128$ MNIST exponent is nearly identical and
$m=64$ is modestly steeper, but the MNIST across-seed standard deviations are
large enough that this difference is not decisive.  At $N=12800$, the
learned/random ratios were $.6408\pm.0129$ and $.6711\pm.0033$, and code nnz
fractions were $.13063\pm.00025$ and $.07421\pm.00014$, confirming learning
relative to the random baseline.

\begin{figure}[H]
  \centering
  \includegraphics[width=0.9\linewidth]{dl_part3_naxis.png}
  \caption{MNIST N-axis fits; triangular markers identify capped runs.}
\end{figure}

\subsection{Atoms}

To visualize an atom $a$ in whitened PCA coordinates, I applied the inverse
linear transform
\[
 a_{\rm pixel}=a\,\mathrm{diag}(\sqrt{\lambda_1},\ldots,
 \sqrt{\lambda_k})V_k^\top.
\]
I did not add the data mean to a direction.
Both dictionaries contain recognizable digit and stroke templates.  The
$m=128$ dictionary contains more writing-style variants and near-duplicate
parts.  Greater overcompleteness lowers absolute held-out loss and makes codes
sparser, while individual atoms remain interpretable.

\begin{figure}[H]
  \centering
  \includegraphics[width=0.49\linewidth]{dl_part3_atoms_m64.png}\hfill
  \includegraphics[width=0.49\linewidth]{dl_part3_atoms_m128.png}
  \caption{Inverse-whitened MNIST atoms for $m=64$ (left) and $m=128$ (right).}
\end{figure}
\FloatBarrier

\section{Part 4 (bonus): spectral bias}

\subsection{Derivation and generator check}

Since $\E[b_kg_k]=0$ and coordinates are independent,
\[
 \E[z]=0,\qquad \E[zz^\top]=\operatorname{diag}(p_1v_1,\ldots,p_mv_m)
 =\operatorname{diag}(\tau_1,\ldots,\tau_m).
\]
Consequently, for $x=A^\top z+\epsilon$,
\[
 \operatorname{Cov}(x)=A^\top\operatorname{diag}(\tau_k)A+\sigma^2I.
\]
If the rows $a_k^\top$ of square $A$ are orthonormal, then
$\operatorname{Cov}(x)a_k=(\tau_k+\sigma^2)a_k$: each atom is an eigenvector
with eigenvalue $\tau_k+\sigma^2$.

Here $p_k=ck^{-\alpha}$ and $v_k=1$, so expected sparsity is
$\sum_{k=1}^m p_k=c\sum_{k=1}^m k^{-\alpha}$, which decreases with $\alpha$
(for all $k>1$).  For $m=16,c=.45$, its values are 2.079, 1.521, and 1.175 at
$\alpha=.75,1,1.25$.  If instead $v_k=1/p_k$, then $\tau_k=1$ and every signal
eigenvalue is equal.  This explains why the variance-normalized Part 0
generator cannot test spectral bias.  The empirical coordinate variances in
Figure~\ref{fig:p4var} track the intended unequal $\tau_k$.

\begin{figure}[H]
  \centering
  \includegraphics[width=0.7\linewidth]{dl_part4_variances.png}
  \caption{Empirical code variances for five seeds at $\alpha=1$.}
  \label{fig:p4var}
\end{figure}

\subsection{Recovery ordering and algorithm dependence}

With alternating minimization, high-energy atoms are generally recovered
first and the separation strengthens as $\alpha$ increases.  Seed variation
is visible but smaller than the overall high-to-low-energy trend.  Unrecovered
atoms are shown above the budget line.

\begin{table}[H]
\centering
\caption{Part 4 recovery summaries.  Mean time is first averaged within each
seed over recovered atoms; ranges pool recovered atom--seed pairs.}
\begin{tabular}{llrrr}
\toprule
Method & $\alpha$ & recovered fraction & mean recovery time & pooled range \\
\midrule
Alt. min. & .75 & $1.000\pm.000$ & $17.28\pm5.26$ & 2--58 \\
Alt. min. & 1.00 & $.975\pm.056$ & $24.11\pm7.37$ & 1--79 \\
Alt. min. & 1.25 & $.900\pm.084$ & $32.95\pm11.69$ & 1--113 \\
MSP       & 1.00 & $1.000\pm.000$ & $2.83\pm.29$ & 1--6 \\
\bottomrule
\end{tabular}
\end{table}

\begin{figure}[H]
  \centering
  \includegraphics[width=0.98\linewidth]{dl_part4_ordering.png}
  \caption{Alternating-minimization recovery times.  Red triangles are atoms
  not recovered within 120 blocks.}
\end{figure}

At $\alpha=1$, MSP recovers all 80 atom--seed pairs within six iterations and
shows little persistent ordering by energy.  Unlike coordinate inference and
a gradual dictionary gradient step in alternating minimization, MSP applies a
fourth-moment stretch followed by a polar projection, updating the entire
orthogonal basis jointly.  This global coupled projection rapidly enters its
local recovery phase and largely removes the slow, coordinate-by-coordinate
spectral ordering.

\begin{figure}[H]
  \centering
  \includegraphics[width=0.82\linewidth]{dl_part4_algorithm.png}
  \caption{Identical $\alpha=1$ data and initializations: alternating
  minimization (left) versus MSP (right).}
\end{figure}

\section{Reproducibility and limitations}

All required dimensions satisfy $D\le64$, $m\le256$, and $N\le20000$.
Random seeds, configurations, stopping states, censored/capped flags, and
per-seed measurements are stored in JSONL logs.  The main limitations are the
short floor-limited transient in Part 1(a), unstable eval-excess extrapolation
in Part 1(b), finite-sample and algorithm dependence of the observed Part 2
transition, and 15 capped MNIST N-axis runs.  These caveats are retained rather
than silently filtering runs.

\section{AI usage disclosure}

I used an OpenAI coding assistant through the Pi coding-agent interface to
help inspect the assignment, audit the experimental plan, implement and review
NumPy experiment code, diagnose the local Python/LAPACK environment, run the
experiments, check plots and logs, and prepare a draft of this report.  The
assistant also helped identify and correct issues including the sign of the
Bernoulli--Gaussian excess kurtosis, the orthogonal reparameterization at
$\theta=1$, inverse-whitening of atoms, stationarity-based stopping, and
continuing Part 2 runs beyond the first-passage time before evaluating loss.
All numerical claims in this document come from the repository's saved logs
and reproducible scripts.  \todo{Before submission, personally verify the
code, figures, derivations, wording, repository URL, and that this disclosure
accurately describes your use and understanding.}

\end{document}
